\documentclass[UTF8]{ctexart}
\usepackage{geometry}
\usepackage{amsmath,amssymb,amsfonts}
\usepackage{array,booktabs}
\geometry{a4paper,margin=2.5cm}

\begin{document}

\begin{center}
\LARGE\bfseries 第五章作业题单
\end{center}

\section*{教材《概率论与数理统计》}

\subsection*{第五章习题（pp.125--126）}

\begin{enumerate}
\renewcommand{\labelenumi}{\textbf{\arabic{enumi}.}}
\setcounter{enumi}{1}
\item 
\begin{enumerate}
\item[(1)] 一保险公司有$10\,000$个汽车投保人，每个投保人索赔金额的数学期望为280美元，标准差为800美元，求索赔总金额超过$2\,700\,000$美元的概率。
\item[(2)] 一公司有50张签约保险单，各张保险单的索赔金额$X_i,\ i=1,2,\cdots,50$（以千美元计）服从韦布尔（Weibull）分布，均值$E(X_i)=5$，方差$D(X_i)=6$，求50张保险单索赔的合计金额大于300的概率（设各保险单索赔金额是相互独立的）。
\end{enumerate}

\setcounter{enumi}{3}
\item 设各零件的质量都是随机变量，它们相互独立，且服从相同的分布，其数学期望为$0.5\ \mathrm{kg}$，均方差为$0.1\ \mathrm{kg}$，问$5\,000$个零件的总质量超过$2\,510\ \mathrm{kg}$的概率是多少？

\setcounter{enumi}{5}
\item 一工人修理一台机器需两个阶段，第一阶段所需时间（以$\mathrm{h}$计）服从均值为0.2的指数分布，第二阶段所需时间服从均值为0.3的指数分布，且与第一阶段独立。现有20台机器需要修理，求他在$8\ \mathrm{h}$内完成的概率。

\setcounter{enumi}{11}
\item 一公寓有200户住户，一户住户拥有汽车辆数$X$的分布律为
\[
\begin{array}{c|ccc}
X & 0 & 1 & 2\\
\hline
p_k & 0.1 & 0.6 & 0.3
\end{array}
\]
问需要多少车位，才能使每辆汽车都具有一个车位的概率至少为0.95？
\end{enumerate}

\section*{教材《概率论及其应用》}

\subsection*{第十章习题（p.201 / 217）}

\begin{enumerate}
\renewcommand{\labelenumi}{\textbf{\arabic{enumi}.}}
\setcounter{enumi}{6}
\item 设$\{X_k\}$是相互独立的具有同分布的随机变量序列，且有均值$\mu$和有限的方差。令$S_n=X_1+\cdots+X_n$，证明大数定律对随机变量序列$\{S_n\}$不成立。但是，如果$na_n\to 0$，则$\{a_nS_n\}$服从大数定律。

提示：计算$\mathrm{Var}(S_1,\cdots,S_n)/n$。

\item 设$\{X_k\}$是随机变量序列，其中$X_k$只依赖于$X_{k-1}$和$X_{k+1}$而与其它的$X_j$相互独立，且$\mathrm{Var}(X_k)$有界，证明$\{X_k\}$服从大数定律。
\end{enumerate}

\end{document}
